\section{Introduction}
\label{sec:introduction}

A conifold transition relates exceptional curves on a small
resolution to vanishing three-spheres on a smoothing.
The relations among those spheres give only the kernel of the
homomorphism sending each sphere to its class in middle homology.
The integral placement of the classes themselves enters the
integral monodromy representation and the saturation of the
vanishing-cycle lattice. For the compact family studied here,
integral topology shows that the lattice is primitive and
determines this placement up to an integral symplectic change
of basis. The coordinates of the spheres in an independently
constructed geometric marking remain to be identified.
We recover those coordinates and the corresponding monodromy
matrices from geometrically represented period observations
on a family of Calabi--Yau threefolds.

The family consists of complete intersections of bidegrees
$(1,1)$, $(1,1)$ and $(2,2)$ in
$\PP^3\times\PP^3$, with polarization $\OO(2,2)$.
Its central fibre has eight ordinary double points, divided
into two groups $I_f,I_g$ of four. The projective small
resolution has second Betti number four; smooth fibres
have second Betti number two and free middle homology
of rank $134$. Orient the local vanishing spheres
$\delta_i$ so that their only integral relations are
\begin{equation}\label{eq:intro-relations}
 \sum_{i\in I_f}\delta_i=0,\qquad
 \sum_{i\in I_g}\delta_i=0.
\end{equation}
These two relations are supplied as geometric data.
The unknown is the marked inclusion of their rank-six
span in middle homology. Its saturation and torsion-free
quotient are determined independently in
Proposition~\ref{prop:primitive-vanishing-lattice};
Corollary~\ref{cor:vanishing-normal-form} gives its symplectic
normal form.

\subsection{The principal result}

Let
\[
 R=\{p\in\CC^8:\ \sum_{i\in I_f}p_i=
                       \sum_{i\in I_g}p_i=0\}.
\]
The relative residue form is denoted by $\Psi$.
An evaluation name for a parameter supplies rational
complex approximations to any requested absolute error.
This is the meaning of access to an arbitrary complex
parameter below; computable parameters give ordinary
algorithms.

\begin{theorem}\label{thm:main-reconstruction}
There is a terminating selection of algebraic central
coefficients and six algebraic coefficient directions
with the following properties.
One can construct a rational-radius ball $B_{\rm ev}\subset R$
and a holomorphic coefficient section defining a flat
projective family $X_p$ with
\[
 \int_{\delta_i}\Psi_p=p_i,\qquad
 B_{\rm ev}^{\circ}
       =B_{\rm ev}\setminus\bigcup_i\{p_i=0\}.
\]
This is exactly its smooth-fibre locus.
The coefficient section, local normal forms and residue
units admit evaluation and bounded derivatives of every
prescribed finite order on specified smaller domains.

For every named $p\in B_{\rm ev}^{\circ}$, a terminating
procedure constructs a geometric symplectic marking
$\gamma_1,\ldots,\gamma_{134}$ and recovers
\[
 D:\ZZ^8\longrightarrow H_3(X_p;\ZZ),\qquad
                        D e_i=\delta_i .
\]
It returns the marked columns, their primitive integral image,
and the individual monodromy matrices, and recovers the
topologically determined quotient $\ZZ^{128}$. These matrices determine
the homology action of every loop in $B_{\rm ev}^{\circ}$.
The observations are the eight local periods and the
$1072$ complex jumps of the marked periods around
represented individual node meridians.
Their absolute enclosures are constructed and refined
until a strict integer-separation test succeeds.
Neither the coordinates of the vanishing spheres nor
the homology coordinates of transported endpoints are
inputs to this forward calculation.
\end{theorem}

The proof occupies Sections~\ref{sec:algebraic-family}--%
\ref{sec:integral-recovery}; its final assembly is
Theorem~\ref{thm:integral-reconstruction}.
There is no a priori bound on the integer entries and
no comparison condition among the nonzero $p_i$.
Termination is pointwise in the smooth parameter.
The result does not give a practical running-time
bound near the discriminant. In particular, we do not
claim numerical values for the selected coefficients or
the attachment matrix in the constructed marking. Its six
nonzero Smith invariants are all one by topology.
The local spheres are globally flat and the marking is
transported from a reference fibre. Thus the matrix is constant
under this transport, and recovery at the reference fibre suffices.
Direct observations at other fibres remain available with the
pointwise error bounds proved below.

The finite inverse step has a simple form. If
$\delta_i=\sum_b(d_i)_b\gamma_b$ and
$J_{ab}=\langle\gamma_a,\gamma_b\rangle$, a positive
individual meridian gives
\begin{equation}\label{eq:intro-period-jump}
 \Delta_{ai}
   =\int_{(F_i)_\#\gamma_a}\Psi_p-\int_{\gamma_a}\Psi_p
   =p_i(Jd_i)_a .
\end{equation}
Thus one holomorphic form suffices for these rank-one
variations; its full period functional need not be
injective on integral homology. Unimodularity of $J$
recovers $d_i$ from the integers $(Jd_i)_a$.
The substantive construction is the forward access
to \eqref{eq:intro-period-jump}, with valid absolute
errors and without knowing those integers.

\subsection{Geometric construction and error control}

We first select an algebraic member of the toric
eight-node construction by finite ideal and
nonvanishing tests. Algebraic jets at the nodes give
the rank-six derivative of its period map and a
transverse coefficient slice. Quantitative implicit
functions, Morse coordinates and an explicitly
integrated residue expansion then give an inverse
period section on a known neighbourhood.
The analytic estimates include bounds on a larger
complex domain, so that differentiation and
integration retain controlled errors.

An algebraic smooth reference fibre provides the
integral marking. Semialgebraic triangulation
constructs its chain complex and integral homology.
Face-compatible smoothing gives actual finite
piecewise smooth representatives, with first and
second derivative bounds. Equation-based horizontal
transport moves these representatives to the target
and around each individual meridian. The first-jet
equation controls the transported tangent vectors,
and quadrature of the pulled-back residue form gives
absolute enclosures for the period differences.

The denominator $p_i$ may be arbitrarily small.
Lemma~\ref{lem:integer-separation} bounds the unknown
integer height from the observed numerator and
denominator enclosures, then tests whether the ratio
error is strictly below $1/2$. Refinement eventually
succeeds at every smooth target.
The same calculation gives a sufficient output
precision depending on the local period and true
integer height. The latter enters the explanatory
bound, not the adaptive input.

\subsection{Applications on the same fibres}

The spectral restriction map and the recovered
attachment map belong to one exact sequence.
Writing $V_{\rm lost}$ for the two-dimensional
quotient of divisor classes by the surviving
classes, the curve restriction map is
\[
 \mathsf r:V_{\rm lost}\longrightarrow\RR^8,\qquad
 \mathsf r(x_f,x_g)=x_fv_f+x_gv_g,\qquad
                    \im\mathsf r=\ker(D\otimes\RR).
\]
For the actual Ricci-flat metric, the weak two-form
energy is the capacity-weighted norm of this map.
The marked image of $D$ contains additional integral
information which that real quadratic form does not
determine.

The analytic results of \cite{W_Hodge,W_HYM} apply after
their central geometric and elliptic hypotheses have been
verified for the selected coefficients. Section~\ref{sec:analytic-family}
performs this verification on a smaller full ball
$B_{\rm an}\subset B_{\rm ev}$, whose radius need not be
computable. Appendix~\ref{app:analytic-inputs} records the
complete statements used in these applications.

Theorem~\ref{thm:integrated} identifies the recovered map
with the map in the attachment sequence and relates the
analytic constructions on these same fibres. The nonflat
rank-four bundles have determinant $\OO$, $c_2H=32$, and
small deformations described by two $2\times2$ matrices,
one in each direction between the two multiplicity spaces.
On the polystable locus they admit Hermitian--Einstein metrics.
The invariant manifold of metric flow on a fixed bundle is
tangent there to the complete seven-dimensional low spectral
subspace, including the kernel. Its finite Hessian is relatively
equivalent to the squared norm of the infinitesimal unitary
action on the matrix parameters. Minimizing the deformation
norm over this action and the complementary spatial gauge
directions gives the harmonic metric on the varying-bundle
moduli space. The proof identifies these maps and quadratic
forms using the estimates of \cite{W_HYM}.

The integral observations use represented middle cycles and
remain independent of the bundle metric. Their paths lie in
$B_{\rm ev}^{\circ}$ and may leave $B_{\rm an}$.
The bounded residue units in $\tau_i=p_i/U_i$ relate the
Hodge capacities to the sufficient absolute precision for
direct recovery at a target fibre. Recovering the constant
integer matrix at the reference fibre avoids this
target-dependent precision requirement.

\subsection{Relation to previous work}

The Picard--Lefschetz formula and its integral
topological interpretation are classical
\cite[Expos\'e~XIV]{Deligne}. Effective recovery of
homological monodromy from continued periods is
already an explicit part of Lairez, Pichon-Pharabod
and Vanhove \cite[\S3.5.2, equation~(38), arXiv-v2
text]{LPPV}. Their work develops effective homology
and periods for projective hypersurfaces and includes
topological results for complete intersections in
projective space. A product-ambient complete
intersection does not become such a complete
intersection merely by the Segre embedding.
We do not claim the rank-one inversion in
\eqref{eq:intro-period-jump} as a new method.

Sert\"oz \cite[\S\S1.3,1.6]{Sertoz} gives rigorous
period continuation for hypersurfaces and distinguishes
transported homology bases from tracking actual
representatives. Our construction supplies represented
cycles and controlled transport for the present
period-coordinate family, including its generally
transcendental coefficient section.
Lairez--Sert\"oz \cite[Theorem~17]{LairezSertoz}
instead give an arithmetic separation theorem for
periods of algebraic quartic surfaces. Here the
quantity to separate is an integer multiple of
one nonzero local period with a measured denominator;
the stopping argument is elementary once its
geometric enclosures have been established.

Vanishing-period coordinates and logarithmic
variation in Calabi--Yau conifold transitions are
studied by Lee, Lin and Wang
\cite[\S3.3 and Proposition~3.16]{LLW}.
We require, in addition, coefficient selection,
a known inverse neighbourhood, a geometric integral
marking and absolute observation errors.
Map-valued semialgebraic triangulation is an
established constructive input
\cite[Theorem~5.43 and Algorithm~11.44]{BPR}.
The smoothing of its chains used here retains
common faces and explicit derivative control;
we do not assume a differentiable triangulation
with those bounds.

The central compact transition and its divisor
data come from \cite{W_Toric,W_HodgeRay}.
The relative Ricci-flat Hodge theorem and the
nonflat bundle, HYM and harmonic quotient results
are those of \cite{W_Hodge,W_HYM}.
The present contribution is the terminating geometric
observation construction and its integral output
on a family to which those analytic conclusions
are proved to apply. The final section records
the actual connecting maps and an example with
independent base and bundle scales.
